\section{Related work}
\label{sec:related}

A work is cited here only if it was read in full. Four have been. This is
enough to position against and nowhere near a survey; \cref{sec:limitations}
records that as the largest known gap rather than disguising it.

\paragraph{Both FEA-agent neighbours route around the documentation problem,
and both say so.}
AbaqusAgent~\cite{abaqusagent2026} does not ingest the Abaqus documentation at
all. It curates 104 worked cases --- 71 from the benchmark manual, 33
textbook-style --- embedded into a single vector store, with a whole case as
the retrieval unit and an explicit prompt instruction that a retrieved case is
a template for structure, not physics. The authors give the reason plainly:
the documentation of a closed-source commercial package is not readily
available, and curating it is a significant challenge under limited resources.
PDE-Agents~\cite{pdeagents2026} likewise holds material properties, failure
patterns and run lineage in its knowledge graph rather than manuals.

\paragraph{The measured benefit of retrieval in this neighbourhood is
efficiency, not accuracy.}
AbaqusAgent's ablation (20 cases, one model, temperature 0) reports 90\%
success with retrieval against 80\% without, and 634{,}690 tokens against
2{,}203{,}778 --- $3.5\times$ fewer tokens and $2.2\times$ less time for
$+10$pp, an accuracy delta they do not claim is significant at that $n$. This
directly shapes our endpoint choice (\cref{sec:evaluation}). Two cautions
carry over. Their evaluation set of 50 draws 40 cases from the 104-case
retrieval library, which inflates any retrieval-on condition. And the same
table run on a different model moves from 90\%/90\% to 65\%/45\% at fixed
architecture: \textbf{model choice moved the result more than the ablated
feature did}, which is why every arm here pins and names its model.

\paragraph{Integration pattern beats index choice --- two independent groups.}
PDE-Agents states it outright: integration pattern, rather than knowledge
content, decides whether graph-augmented retrieval helps or hinders.
\cite{grepallyouneed2026} finds grep beating vector retrieval in its first
experiment, but its own conclusion is that the harness and tool-calling style
dominate the retrieval strategy. If that holds, ablating \emph{how} a
capability is wired into the loop is more informative than ablating
\emph{which} retriever is used --- which is the shape our evaluation already
has.

\paragraph{What does not transfer, recorded because it was nearly over-read.}
\cite{grepallyouneed2026} benchmarks on LongMemEval, conversational memory
where answers are often licensed by a small set of distinctive strings: the
best possible case for lexical matching and a poor analogy for technical
documentation. \cite{chunking2026} reports no measurable benefit from chunk
overlap, for a configuration (open-domain Wikipedia QA, a sparse-learned
retriever, a small generator) that shares none of our parameters. Both findings
are real and neither bears on our setting.

\paragraph{The opening, stated carefully.}
Nobody in this neighbourhood has made a full documentation corpus traversable;
both neighbours hand-curate a small structured knowledge base instead. Two
risks before claiming that gap. It is adjacent to graph-augmented retrieval,
which PDE-Agents has already applied in this exact domain, so the difference
must be principled --- a hierarchy \emph{derived} from document structure
versus a \emph{curated} graph of facts. And given the integration-beats-index
finding, a traversable tree without a traversal policy would likely
underperform its own promise.
